\chapter*{阅读指引}

分类器在一批训练邮件上答对很多次，仍可能只是记住了这批邮件的特征与标签；即使当前测试错误很少，换一批训练邮件也可能学出另一条规则。理论部分要回答的是：有限样本在什么条件下足以支持对未来风险的判断。要回答它，必须同时说清观测怎样产生、允许哪些决策、如何计算损失、算法能看到什么，以及“成功”是相对于哪个基准、在哪些随机变化下成立。

基础三章已经区分了任务、模型、损失和反馈，数学两卷提供了组合维度、集中不等式、随机复杂度、信息量、分布距离和优化条件等工具。本部分把这些工具接到具体的学习问题上，追问三件事：候选规则的范围怎样受控，样本怎样影响算法的选择，在样本上算出的评价怎样接近总体风险。涉及神经网络时，只建立推导所需的形式与条件；结构设计、反向传播和训练实践见\BookVolumeRef{vol:neural-network-models}{《模型：神经网络模型》}。

六章大致按“每次改变一类条件”的顺序安排。《机器学习理论基础》用邮件二分类说明学习过程，给出PAC可学习的正式定义，说明归纳偏置为什么必要，把风险分解为贝叶斯风险、近似误差与类内超额风险，并建立各章共用的风险比较和可计算学习的定义。《二分类学习的可学习性》由此出发，从有限个候选扩展到无限假设空间，用生长函数与VC维刻画可学习性，再核对统计保证能否由程序在可承受的资源内实现。《多分类学习的可学习性》把标签从两类扩大到多类，说明有限标签时Natarajan维仍统一刻画学习，一般标签下则由图维刻画一致收敛、由DS维刻画可学习性。《一般学习问题的可学习性》把零一损失换成一般损失，先建立数据依赖的复杂度、样本压缩、算法稳定性和信息依赖四类分析工具，再用它们分析模型选择、集成学习、正则化和梯度下降的风险保证。《深度学习的可学习性》研究网络训练怎样选解、宽网络何时近似固定核或移动特征，再给出网络的风险保证，最后用这些分析审视双重下降、延迟泛化和神经坍缩三个现象。《迁移学习和持续学习的可学习性》最后放开同分布和数据可重访的前提，研究源域与目标域的联系、已有知识是否充分、历史信息怎样保留，才能控制目标风险和历史风险。

六章共用同一套风险、概率量词和可测性约定，又逐步改变标签、损失、选解方式及训练样本与评价样本的生成规律。因此，每次复用前文结论，都要先核对这些对象和条件是否仍然相同。所得保证约束的是给定问题中的风险或超额风险；损失是否覆盖实际后果、比较基准是否可接受、保证能否推广到允许范围之外的环境，本部分不作判断。下一部分将把这些使用问题逐项纳入评价与处置。
